% Image credits for the photographs and AI illustrations of Book 4 (University, Year 2).
% Machine-readable license records live in images/book4/CREDITS.md; keep the
% two files in sync. The Book 4 writer owns this file: add one \item per
% photograph (inside the itemize, by chapter) as photographs land.
% Portuguese (Brazilian) edition: translated from frontmatter/image-credits-book4.tex;
% every credit, name, dataset and licence kept exactly.
\chapter*{Créditos das imagens}

Os desenhos, os mapas e os gráficos deste livro são originais; os mapas e os
gráficos desenhados a partir de dados publicados citam sua fonte na legenda ou
nesta lista. As fotografias são usadas sob suas respectivas licenças livres, com
agradecimentos a seus autores; os links completos para cada fotografia e as URLs
das licenças estão listados em \texttt{images/book4/CREDITS.md}, no repositório
do livro.

\bigskip
Ao lado das fotografias e dos desenhos originais, algumas figuras são ilustrações
geradas por inteligência artificial, produzidas com a geração de imagens da OpenAI
a partir de prompts escritos pelos editores do livro e revisadas quanto à precisão
científica antes de sua inclusão. Elas mostram apenas cenas genéricas, nunca um
lugar real, uma pessoa real ou um espécime real. O prompt completo de cada imagem
gerada está registrado em \texttt{images/book4/ai/PROMPTS.md}, no repositório.

\bigskip
\noindent Fotografias e dados, por capítulo:
\begin{itemize}\small
  \item Cap.~1 (o meio interestelar e a formação das estrelas): o disco protoplanetário de HL Tauri, ALMA (ESO/NAOJ/NRAO), CC BY 4.0; os Pilares da Criação na Nebulosa da Águia, NASA, ESA e Hubble Heritage Team (STScI/AURA), domínio público.
  \item Cap.~2 (formação dos planetas e meteoritos): o condrito carbonáceo de Allende e o padrão de Widmanstätten do meteorito metálico de Gibeon, James St.~John, CC BY 2.0.
  \item Cap.~3 (os elementos na Terra): xenólito mantélico de lherzolito em basalto, Kilbourne Hole, James St.~John, CC BY 2.0.
  \item Cap.~5 (os magmas e seu alojamento): dique perto de Shiprock, Dexter Perkins, CC0 (Wikimedia Commons).
  \item Cap.~6 (as erupções vulcânicas): fonte de lava no cone Pu`u~`\=O`\=o do Kīlauea, J.\,D. Griggs, USGS, domínio público; a coluna eruptiva do Monte Santa Helena, 18 de maio de 1980, Austin Post, USGS, domínio público; a coluna eruptiva do Monte Pinatubo, 12 de junho de 1991, Dave Harlow, USGS, domínio público.
  \item Cap.~7 (a fusão do manto): amostra de mão de dunito, Mike Davis (GeoDIL), CC0.
  \item Cap.~8 (a diferenciação magmática): camadas de cromitito do Complexo de Bushveld, Kevin Walsh, CC BY 2.0; gabros acamadados da intrusão de Skaergaard, C. Tegner, CC BY 4.0.
  \item Cap.~10 (reações e fácies metamórficas): xisto azul com lawsonita, Ward Creek, Califórnia, e eclogito, Almenningen, Noruega, James St.~John, CC BY 2.0.
  \item Cap.~11 (termobarometria e trajetórias P--T--t): xisto com granada, James St.~John, CC BY 2.0; grão de coesita em eclogito, entre polarizadores cruzados, J.\,R. Smyth, CC BY-SA 1.0.
  \item Cap.~12 (a fusão da crosta: migmatitos e granitos): afloramento de migmatito em Nashua, New Hampshire, Kevin M. Gill, CC BY 4.0; veio de leucogranito, Anne Burgess (Geograph), CC BY-SA 2.0.
  \item Cap.~13 (anatomia de um cinturão de montanhas: os Alpes): lavas em almofada do ofiolito do Chenaillet, Lionel Duchoiselle, CC BY-SA 3.0; o Matterhorn visto de Zermatt, Ximonic (Simo Räsänen), CC BY-SA 4.0.
  \item Cap.~14 (cinturões de colisão e cinturões andinos comparados): o Planalto Tibetano entre o Nepal e Shigatse, Jan Kranendonk, CC BY-SA 4.0; o Altiplano boliviano, Luca Galuzzi, CC BY-SA 2.5.
  \item Cap.~15 (maciços antigos e o ciclo de Wilson): Siccar Point, falésia sudeste, dave souza, CC BY-SA 4.0.
  \item Cap.~16 (bacias sedimentares comparadas): plataforma de perfuração semissubmersível no mar do Norte, Erik Christensen, CC BY-SA 3.0.
  \item Cap.~17 (ilhas oceânicas): o Atol Kure visto da Estação Espacial Internacional, NASA, domínio público.
  \item Cap.~18 (a Terra pré-cambriana): komatiito spinifex em Pyke Hill, Ontário, James St.~John, CC BY 2.0; carbonato de capa da Formação Nuccaleena, Austrália Meridional, James St.~John, CC BY 2.0.
  \item Cap.~19 (a Terra fanerozoica): pilares de creta de Monument Rocks, Kansas, James St.~John, CC BY 2.0; o nível Bonarelli, Picenus, CC BY-SA 4.0.
  \item Cap.~20 (fósseis e paleoambientes): \emph{Waptia} e \emph{Marrella} do Folhelho Burgess, James St.~John, CC BY 2.0; o espécime de Berlim do \emph{Archaeopteryx}, H.~Raab, CC BY-SA 3.0; \emph{Zoophycos} em Sagres, Amada44, CC BY 4.0.
  \item Cap.~21 (termodinâmica da atmosfera): bigorna de uma tempestade supercelular, Stephen Corfidi, NOAA, domínio público; sondagens de radiossonda do Integrated Global Radiosonde Archive, I. Durre e outros, NOAA NCEI, domínio público.
  \item Cap.~22 (sistemas meteorológicos): um sistema de baixa pressão sobre a Islândia, NASA/GSFC, MODIS Rapid Response Team, Jacques Descloitres, domínio público; o olho do furacão Isabel visto da Estação Espacial Internacional, Ed Lu, NASA, domínio público.
  \item Cap.~23 (variabilidade climática): índices da Oscilação do Atlântico Norte e Oceanic Niño Index, NOAA Climate Prediction Center, domínio público.
  \item Cap.~24 (arquivos e indicadores climáticos): camadas anuais do testemunho de gelo GISP2 a 1855~m, Anthony Gow, U.S. Army Corps of Engineers CRREL (NOAA Paleoclimatology), domínio público; seção transversal de um espeleotema de travertino, Diamond Caverns, Kentucky, James St.~John, CC BY 2.0; o empilhamento bentônico LR04 de L.\,E. Lisiecki e M.\,E. Raymo (2005), o registro de temperatura EPICA Dome C de J. Jouzel e outros (2007) e a composição antártica de \ce{CO2} de B. Bereiter e outros (2015), NOAA NCEI Paleoclimatology, domínio público.
  \item Cap.~25 (o forçamento orbital e as eras glaciais): elementos orbitais e insolação de Berger e Loutre (1991), NOAA NCEI Paleoclimatology, domínio público.
  \item Cap.~26 (o clima ao longo do tempo geológico): registro isotópico bentônico cenozoico CENOGRID, T. Westerhold e outros (2020), PANGAEA, CC BY 4.0; compilação fanerozoica de \ce{CO2} de G.\,L. Foster, D.\,L. Royer e D.\,J. Lunt (2017), \emph{Nature Communications}, CC BY 4.0.
  \item Cap.~27 (a mudança climática antrópica): Sexto Relatório de Avaliação do IPCC, Grupo de Trabalho~I, dados estendidos do Anexo~III, CC BY 4.0.
  \item Cap.~28 (ameaças e riscos geológicos): depósitos de fluxo de detritos em Caraballeda, Venezuela, Lawson Smith, U.S. Army Corps of Engineers, domínio público; uma mesquita de pé após o tsunami de 2004 em Aceh, USGS, domínio público.
  \item Cap.~29 (jazidas minerais): a mina de Bingham Canyon vista da Estação Espacial Internacional, NASA (tripulação da Expedição~15), domínio público; gossan da mina de Flambeau, Wisconsin, James St.~John, CC BY 2.0; dados de teor e tonelagem de D.\,A. Singer, V.\,I. Berger e B.\,C. Moring (2008), USGS, domínio público.
  \item Cap.~30 (geologia da energia): usina geotérmica perto de Geyserville, Califórnia, Departamento de Defesa dos EUA, domínio público; densidades do dióxido de carbono do NIST Chemistry WebBook (Standard Reference Database 69).
  \item Cap.~31 (métodos de campo e de laboratório II): augita poiquilítica englobando plagioclásio, Strekeisen, CC BY-SA 4.0; lâmina delgada de um granada--mica xisto, Jackdann88, CC BY 3.0; microssonda eletrônica do Serviço Geológico da Finlândia, Kallerna, CC BY-SA 4.0; resultados de sílica de materiais de referência de rochas, J.\,E. Taggart e D.\,F. Siems (2002), USGS, domínio público.
\end{itemize}
